\chapter{Sensitivity analysis}\label{chap:sensitivity}

Elasticities are rarely estimated for the country and period of a SAM, and
results can depend on them. \manref{cge_sensitivity} shows how much: it
recalibrates the model with elasticities drawn around the values used, runs
the scenario and solves, draw after draw.

\section{Draws}

Each elasticity $\sigma$ of the types chosen (value added, Armington, CET and
exports from composite supply, and the top nest) is multiplied, in draw $k$, by
a factor $u^{(k)}$ around one:
\begin{equation*}
\sigma^{(k)} = u^{(k)}\,\sigma, \qquad u^{(k)} = 1 + s\,v^{(k)},
\end{equation*}
where $s < 1$ is the spread and $v^{(k)}$ is drawn from the uniform distribution
on $[-1, 1]$, $v = 2U - 1$, or from the symmetric triangular distribution on the
same interval, $v = U_1 + U_2 - 1$, with $U$, $U_1$ and $U_2$ independent and
uniform on $[0, 1]$. By default each sector's elasticity has its own factor;
with \cmd{common} one factor applies to every sector for each type of
elasticity. A Leontief top nest (elasticity 0) stays Leontief, and the
elasticities of nests of value added keep the values given. Draw 0, with the
central elasticities, is always added.

For each draw the model is recalibrated from the SAM with every other
calibration setting kept (factor quantities, LES demand, the health link), the
closure is restored, the scenario program applies its shocks and the model is
solved. The model in memory, with its simulations, is restored at the end.

\section{Summary}

For every report measure of chapter~\ref{chap:reporting}, the summary gives the
result with the central elasticities, and, across the $n$ draws that
converged, the mean, the standard deviation and the 5th and 95th percentiles,
taken as order statistics of the sorted results $x_{(1)} \le \dots \le
x_{(n)}$:
\begin{equation*}
p_{5} = x_{(\max(1, \lceil 0.05\,n \rceil))}, \qquad
p_{95} = x_{(\max(1, \lceil 0.95\,n \rceil))}.
\end{equation*}

\section{Interpretation}

The draws do not estimate uncertainty about the elasticities. They show how
the results change over a chosen range around the values used. The range, the
distribution and whether sectors move together should be stated with the
results. Standard errors and bounds of published estimates
(chapter~\ref{chap:elasticities}) can guide the choice of the range.
